% Bagean: metodhe
% Bab: bukti
% Pérangan: aku-ora-bisa

\documentclass[../../../include/open-logic-section]{subfiles}

\begin{document}

\olfileid{mth}{prf}{cnt}

\olsection{Aku Ora Bisa!{}}

Kita kabèh tau tekan wektu nalika rumangsa arep nyerah.
Nanging panjenengan \emph{bisa} nglakoni. Dhosèn lan
asistèn pangajaran, uga kanca-kanca kuliah, bisa
mbantu. Nyuwuna pitulungan marang wong-wong iku!{}
Ing ngisor iki ana sawatara saran kanggo nyegah
rasa kewalahan lan kanggo tumindak nalika
panjenengan rumangsa arep nyerah.

Supaya bisa ngrampungaké soal kanthi becik,
lakonana iki:
\begin{enumerate}
\item \emph{Miwitia sedini mungkin.} Sajroning semèster
  kita padha sibuk lan akèh kang angel ngindhari
  pakulinan nundha-nundha. Salah siji cara paling
  migunani yaiku miwiti tugas omah luwih awal.
  Yèn banjur macet, panjenengan isih nduwèni
  wektu kanggo nggoleki dalan metu saliyané nangis.
\item \emph{Rembugana karo kanca sakelas}. Panjenengan
  ora dhèwèkan. Kanca liya bisa uga nemoni
  kangèlan, nanging ing bab kang béda. Rembugan
  bisa mènèhi sudut pandang anyar marang soal lan
  mbukak dalan kanggo nemokaké jawaban.
  Aja mung nyalin jawabané: nyuwuna pituduh,
  utawa jlèntrèhna ing endi panjenengan macet
  banjur takokna langkah sabanjuré. Yèn wis
  ngerti, tulungana uga wong liya. Mbantu lan
  nerangaké marang kanca uga bakal njalari
  pangertèn panjenengan dhéwé luwih jero.
\item \emph{Nyuwuna pitulungan.} Ana akèh sumber
  pitulungan kang kasedhiya. Dhosèn lan asistèn
  pangajaran ana kanggo mbantu lan \emph{péengin}
  panjenengan kasil. Wong-wong mau bisa mbantu
  nggarap soal lan nemtokaké ing langkah endi
  panjenengan ngalami kangèlan.
\item \emph{Ngaso sedhéla.} Yèn panjenengan macet,
  \emph{mbokmenawa} amarga wis kakehan wektu nyawang
  soal kang padha. Ngasoa sedhéla, ngombé téh,
  utawa garapen soal liya sawetara wektu;
  banjur balia kanthi pikiran kang seger.
  Coba dipikir maneh sawisé turu.
\end{enumerate}

Sadarana yèn kabèh cara iki mbutuhaké
panjenengan wis miwiti bukti adoh sadurungé
wektu nglumpukaké. Yèn lagi miwiti jam loro
esuk ing dina sadurungé tugas diklumpukaké,
saran-saran mau mbokmenawa ora pati mbantu.

Iki mbokmenawa krungu kaya pangeling kang
peteng, nanging ngrampungaké bukti iku
tantangan kang pungkasane menehi asil.
Ana wong kang ndadèkaké iki pakaryan,
dadi mesthi ana bab kang bisa disenengi.
Kaya meh kabèh katrampilan, ngrampungaké
soal lan nyusun bukti mbutuhaké latihan.
Panjenengan bisa ndeleng kanca kang katon
gampang nindakaké: mbokmenawa wong mau
wis akèh latihan. Aja gampang nyerah.

Yèn wektu utawa kasabaran kanggo sawijining
soal pancèn entèk, ora apa-apa. Iki ora ateges
panjenengan bodho utawa ora bakal bisa ngerti.
Takokna marang dhosèn utawa kanca carané
ngrampungaké, banjur temokna ing endi
panjenengan salah utawa macet supaya ora
mbalèni kangèlan padha ing mangsa ngarep.
Sawisé iku, cobanen ngrampungaké maneh
tanpa ndeleng jawaban. Ing wektu sabanjuré,
miwitia luwih awal lan nyuwuna pitulungan
luwih awal uga.

\end{document}
